\label{sec:absolute-value-as-distance}

\begin{topicbox}{Absolute Value Is Distance to Zero}
\begin{exposition}
The metric $d$ specializes to absolute value when one point is the origin. This is the geometric reading of $|a|$: how far $a$ sits from $0$ on the line.
\end{exposition}
\end{topicbox}
\begin{propositionbox}{Proposition (Absolute Value Is Distance to the Origin)}
\begin{proposition}[Absolute Value Is Distance to the Origin]
\label{prop:abs-value-is-distance-to-zero}
For every $a\in\mathbb{R}$, $\;|a|=d(a,0)$.
\hyperref[prf:abs-value-is-distance-to-zero]{Go to proof.}
\end{proposition}
\end{propositionbox}
\begin{remark*}[Standard quantified statement]
\[
\forall a\in\mathbb{R}\;(|a|=d(a,0)).
\]
\end{remark*}
\begin{remark*}[Interpretation]
Absolute value is not a separate idea from distance; it is the distance to the origin. The two-sided bound $|x|\le r \Leftrightarrow -r\le x\le r$ is then exactly the statement that $x$ lies within distance $r$ of $0$, i.e.\ in the interval $[-r,r]$.
\end{remark*}
\begin{dependencies}
\begin{itemize}
  \item \hyperref[def:distance-on-real-line]{Distance on the Real Line}
\end{itemize}
\end{dependencies}
